\documentclass[12pt]{article}
\usepackage[utf8]{inputenc}
\usepackage[T1]{fontenc}
\usepackage{hyperref}
\title{Wulna Language Resources for K-12 Curriculum Development}
\author{Compiled by Claude Sonnet 5.5 (claude-sonnet-5-5) for Andr\'as Kornai}
\date{}
\begin{document}
\maketitle

\section{Introduction}
Wulna (also Wuna; ISO 639-3 wux, Glottocode wuln1239) is an extinct non-Pama-Nyungan language of the Adelaide River region, Northern Territory, Australia. It is tentatively linked with Limilngan in a Limilngan-Wulna family. Jack Wandi, recorded by Gavan Breen at Humpty Doo in 1980--81, was the last known speaker. The people are also called Djerimanga. See \url{https://en.wikipedia.org/wiki/Wulna_language} and \url{https://en.wikipedia.org/wiki/Djerimanga}.

\section{Primary Educational Resources}
No educational materials or revitalization programme were found in these searches.

\section{Linguistic and Reference Materials}
The AIATSIS Austlang entry (code N29) records Breen's 1980--81 recordings with Jack Wandi (MS 1720): \url{https://collection.aiatsis.gov.au/austlang/language/N29?code=N29}. The earliest source is J.\,W.\,O. Bennett's 1869 vocabulary, held by the State Library of New South Wales and annotated by Paul Foelsche; no URL for it was seen. Wulna (N29) has sometimes been confused with Wulwulam (N32) in the literature. Further context: \url{https://en.wikipedia.org/wiki/Darwin_Region_languages} and \url{https://grambank.clld.org/languages/nucl1327}.

\section{Religious and Community Resources}
None were identified in these searches.

\section{Digital and Audio Resources}
Breen's recordings are reported to be available through Paradisec. No Paradisec URL was seen, so none is linked.

\section{Curriculum Development Context}
Material is very sparse: a short 1869 vocabulary and one speaker's 1980--81 recordings. Any heritage curriculum would have to rely on these archival sources. It would need to be developed with descendant communities and the relevant Aboriginal language bodies, whom the searches did not identify. Content would be mostly cultural and historical awareness, with a small amount of recovered vocabulary.

\end{document}